\textbf{Phần 1. Trắc nghiệm lựa chọn}
\setcounterpageref{ex}{0}

\begin{ex} %[5D12N73]
Công thức tính thời gian $t$ trong chuyển động đều là
\choice
{$t = s \times v$}
{$t = v : s$}
{$t = s - v$}
{\True $t = s : v$}
\loigiai{
$t = s : v$.\\
\textbf{Đáp án của câu hỏi là D}
}
\end{ex}

\begin{ex} %[5D12N73]
Một người đi xe máy với vận tốc $40\text{ km/h}$ trong $2{,}5\text{ giờ}$. Quãng đường người đó đi được là
\choice
{$80\text{ km}$}
{$120\text{ km}$}
{\True $100\text{ km}$}
{$90\text{ km}$}
\loigiai{
$s = 40 \times 2{,}5 = 100\text{ (km)}$.\\
\textbf{Đáp án của câu hỏi là C}
}
\end{ex}

\begin{ex} %[5D12H73]
Một ô tô đi quãng đường $126\text{ km}$ hết $2$ giờ $48$ phút. Vận tốc của ô tô là
\choice
{\True $45\text{ km/h}$}
{$42\text{ km/h}$}
{$48\text{ km/h}$}
{$50\text{ km/h}$}
\loigiai{
$2$ giờ $48$ phút $= 2{,}8$ giờ.\\
$v = 126 : 2{,}8 = 45\text{ (km/h)}$.\\
\textbf{Đáp án của câu hỏi là A}
}
\end{ex}

\begin{ex} %[5D12H73]
Một người đi xe đạp với vận tốc $12\text{ km/h}$ đi quãng đường $18\text{ km}$. Thời gian người đó đi là
\choice
{$1$ giờ $20$ phút}
{$1$ giờ $45$ phút}
{\True $1$ giờ $30$ phút}
{$2$ giờ}
\loigiai{
$t = 18 : 12 = 1{,}5\text{ giờ} = 1$ giờ $30$ phút.\\
\textbf{Đáp án của câu hỏi là C}
}
\end{ex}

\begin{ex} %[5D12H73]
Hai địa điểm A và B cách nhau $150\text{ km}$. Hai ô tô cùng khởi hành từ A và B đi ngược chiều nhau với vận tốc lần lượt là $42\text{ km/h}$ và $58\text{ km/h}$. Sau bao lâu hai ô tô gặp nhau?
\choice
{\True $1{,}5\text{ giờ}$}
{$1\text{ giờ}$}
{$2\text{ giờ}$}
{$2{,}5\text{ giờ}$}
\loigiai{
Tổng vận tốc: $42 + 58 = 100\text{ (km/h)}$.\\
Thời gian gặp nhau: $150 : 100 = 1{,}5\text{ (giờ)}$.\\
\textbf{Đáp án của câu hỏi là A}
}
\end{ex}

\begin{ex} %[5D12H73]
Một người đi bộ với vận tốc $4{,}5\text{ km/h}$ trong $40$ phút. Quãng đường đi được là
\choice
{$2\text{ km}$}
{\True $3\text{ km}$}
{$4\text{ km}$}
{$3{,}5\text{ km}$}
\loigiai{
$40$ phút $= \dfrac{2}{3}$ giờ.\\
$s = 4{,}5 \times \dfrac{2}{3} = 3\text{ (km)}$.\\
\textbf{Đáp án của câu hỏi là B}
}
\end{ex}

\begin{ex} %[5D12V73]
Một xe máy khởi hành từ A lúc $7$ giờ $15$ phút với vận tốc $32\text{ km/h}$. Đến $8$ giờ, một ô tô xuất phát từ A đuổi theo với vận tốc $48\text{ km/h}$. Ô tô đuổi kịp xe máy lúc mấy giờ?
\choice
{$9$ giờ $15$ phút}
{$9$ giờ $45$ phút}
{$10$ giờ}
{\True $9$ giờ $30$ phút}
\loigiai{
Từ $7$h $15$p đến $8$h là $45$ phút $= 0{,}75$ giờ.\\
Quãng đường xe máy đã đi: $32 \times 0{,}75 = 24\text{ (km)}$.\\
Thời gian đuổi kịp: $24 : (48 - 32) = 1{,}5\text{ giờ} = 1$ giờ $30$ phút.\\
Thời điểm gặp: $8$ giờ $+ 1$ giờ $30$ phút $= 9$ giờ $30$ phút.\\
\textbf{Đáp án của câu hỏi là D}
}
\end{ex}

\begin{ex} %[5D12V73]
Một thuyền máy đi xuôi dòng một khúc sông từ bến A đến bến B hết $1$ giờ $45$ phút. Vận tốc thực của thuyền là $22\text{ km/h}$, vận tốc dòng nước là $2\text{ km/h}$. Độ dài khúc sông AB là
\choice
{$38\text{ km}$}
{\True $42\text{ km}$}
{$44\text{ km}$}
{$40\text{ km}$}
\loigiai{
Vận tốc xuôi dòng: $22 + 2 = 24\text{ (km/h)}$.\\
$1$ giờ $45$ phút $= 1{,}75$ giờ.\\
$s = 24 \times 1{,}75 = 42\text{ (km)}$.\\
\textbf{Đáp án của câu hỏi là B}
}
\end{ex}

\begin{ex} %[5D12V73]
Một người đi xe máy từ A đến B với vận tốc $36\text{ km/h}$, rồi từ B quay về A với vận tốc $30\text{ km/h}$. Thời gian lúc về nhiều hơn lúc đi là $20$ phút. Quãng đường AB dài là
\choice
{$50\text{ km}$}
{$70\text{ km}$}
{$80\text{ km}$}
{\True $60\text{ km}$}
\loigiai{
Tỉ số vận tốc đi và về: $\dfrac{36}{30} = \dfrac{6}{5} \Rightarrow$ Tỉ số thời gian là $\dfrac{5}{6}$.\\
Hiệu thời gian: $20$ phút $= \dfrac{1}{3}$ giờ.\\
Thời gian đi: $\dfrac{1}{3} : (6 - 5) \times 5 = \dfrac{5}{3}$ giờ.\\
Quãng đường AB: $36 \times \dfrac{5}{3} = 60\text{ (km)}$.\\
\textbf{Đáp án của câu hỏi là D}
}
\end{ex}

\begin{ex} %[5D12V73]
Một đoàn tàu hỏa dài $200\text{ m}$ chạy qua một đường hầm dài $800\text{ m}$ với vận tốc $72\text{ km/h}$. Thời gian từ lúc đầu tàu vào hầm đến lúc đuôi tàu ra khỏi hầm là
\choice
{$40\text{ giây}$}
{$60\text{ giây}$}
{\True $50\text{ giây}$}
{$75\text{ giây}$}
\loigiai{
Tổng quãng đường tàu đi: $800 + 200 = 1,000\text{ (m)}$.\\
Vận tốc: $72\text{ km/h} = 20\text{ m/s}$.\\
Thời gian: $1,000 : 20 = 50\text{ (giây)}$.\\
\textbf{Đáp án của câu hỏi là C}
}
\end{ex}

\begin{ex} %[5D12V73]
Bác Bình đi bộ thể dục $3$ vòng quanh một công viên hình chữ nhật có chiều dài $250\text{ m}$, chiều rộng $150\text{ m}$ với vận tốc $4\text{ km/h}$. Thời gian bác Bình đi hết 3 vòng là
\choice
{$30$ phút}
{\True $36$ phút}
{$40$ phút}
{$45$ phút}
\loigiai{
Chu vi công viên: $(250 + 150) \times 2 = 800\text{ (m)}$.\\
Quãng đường 3 vòng: $800 \times 3 = 2,400\text{ (m)} = 2{,}4\text{ km}$.\\
Thời gian: $2{,}4 : 4 = 0{,}6\text{ giờ} = 36$ phút.\\
\textbf{Đáp án của câu hỏi là B}
}
\end{ex}

\begin{ex} %[5D12C73]
Hai thành phố A và B cách nhau $180\text{ km}$. Cùng lúc $7$ giờ, một ô tô đi từ A đến B và một xe máy đi từ B về A. Hai xe gặp nhau lúc $9$ giờ $15$ phút. Biết vận tốc ô tô gấp đôi vận tốc xe máy. Vận tốc của ô tô là
\choice
{\True $53{,}33\text{ km/h}$ (hay $\dfrac{160}{3}\text{ km/h}$)}
{$50\text{ km/h}$}
{$60\text{ km/h}$}
{$70\text{ km/h}$}
\loigiai{
Thời gian gặp nhau: $9$ giờ $15$ phút $- 7$ giờ $= 2$ giờ $15$ phút $= 2{,}25$ giờ.\\
Tổng vận tốc: $180 : 2{,}25 = 80\text{ (km/h)}$.\\
Vận tốc ô tô: $80 : (2 + 1) \times 2 = \dfrac{160}{3} \approx 53{,}33\text{ (km/h)}$.\\
\textbf{Đáp án của câu hỏi là A}
}
\end{ex}

\textbf{Phần 2. Trắc nghiệm đúng sai}
\setcounterpageref{ex}{0}

\begin{ex} %[5D12H73]
Hai địa điểm A và B cách nhau $135\text{ km}$. Xe máy đi từ A đến B với vận tốc $36\text{ km/h}$, ô tô đi từ B về A với vận tốc $54\text{ km/h}$. Hai xe cùng khởi hành lúc $8$ giờ. Xét tính đúng sai:
\choiceTF
{\True Tổng vận tốc của hai xe là $90\text{ km/h}$.}
{\True Hai xe gặp nhau sau $1$ giờ $30$ phút.}
{Hai xe gặp nhau lúc $9$ giờ $15$ phút.}
{\True Nơi gặp nhau cách A quãng đường $54\text{ km}$.}
\loigiai{
\begin{itemize}
    \item Tổng vận tốc: $36 + 54 = 90\text{ km/h}$. Đúng.
    \item Thời gian gặp: $135 : 90 = 1{,}5\text{ giờ} = 1$ giờ $30$ phút. Đúng.
    \item Gặp nhau lúc: $8$h $+ 1$h $30$p $= 9$h $30$p $\neq 9$h $15$p. Sai.
    \item Cách A: $36 \times 1{,}5 = 54\text{ km}$. Đúng.
\end{itemize}
\begin{itemize}
    \item a) \textbf{Đúng}.
    \item b) \textbf{Đúng}.
    \item c) \textbf{Sai}.
    \item d) \textbf{Đúng}.
\end{itemize}
}
\end{ex}

\begin{ex} %[5D12V73]
Một ca nô đi xuôi dòng khúc sông từ A đến B hết $2$ giờ, đi ngược dòng từ B về A hết $3$ giờ. Vận tốc dòng nước là $3\text{ km/h}$. Xét tính đúng sai:
\choiceTF
{\True Tỉ số vận tốc xuôi dòng và ngược dòng là $\dfrac{3}{2}$.}
{\True Hiệu vận tốc xuôi dòng và ngược dòng là $6\text{ km/h}$.}
{Vận tốc ca nô xuôi dòng là $15\text{ km/h}$.}
{\True Độ dài khúc sông AB là $36\text{ km}$.}
\loigiai{
\begin{itemize}
    \item Tỉ số vận tốc: $3 : 2 = \dfrac{3}{2}$. Đúng.
    \item Hiệu vận tốc: $2 \times 3 = 6\text{ km/h}$. Đúng.
    \item Vận tốc xuôi dòng: $6 : (3 - 2) \times 3 = 18\text{ km/h} \neq 15\text{ km/h}$. Sai.
    \item Khúc sông AB: $18 \times 2 = 36\text{ km}$. Đúng.
\end{itemize}
\begin{itemize}
    \item a) \textbf{Đúng}.
    \item b) \textbf{Đúng}.
    \item c) \textbf{Sai}.
    \item d) \textbf{Đúng}.
\end{itemize}
}
\end{ex}

\textbf{Phần 3. Trả lời ngắn}
\setcounterpageref{ex}{0}

\begin{ex} %[5D12H73]
Một người đi xe máy với vận tốc $38\text{ km/h}$ trong $2{,}5$ giờ. Tính quãng đường theo ki-lô-mét.
\shortans{95}
\loigiai{
\textbf{Đáp số:} $95$\\
$38 \times 2{,}5 = 95\text{ km}$.
}
\end{ex}

\begin{ex} %[5D12H73]
Một ô tô chạy quãng đường $160\text{ km}$ với vận tốc $50\text{ km/h}$. Tính thời gian ô tô đi theo giờ (dùng dấu phẩy nếu có).
\shortans{3,2}
\loigiai{
\textbf{Đáp số:} $3{,}2$\\
$160 : 50 = 3{,}2\text{ giờ}$.
}
\end{ex}

\begin{ex} %[5D12V73]
Hai người đi bộ khởi hành cùng lúc đi ngược chiều nhau trên con đường dài $9\text{ km}$ với vận tốc $4\text{ km/h}$ và $5\text{ km/h}$. Hỏi sau bao nhiêu phút hai người gặp nhau?
\shortans{60}
\loigiai{
\textbf{Đáp số:} $60$\\
$9 : (4 + 5) = 1\text{ giờ} = 60\text{ phút}$.
}
\end{ex}

\begin{ex} %[5D12V73]
Một người đi xe đạp từ nhà đến cơ quan dài $6\text{ km}$ hết $24$ phút. Tính vận tốc của người đi xe đạp theo $\text{km/h}$.
\shortans{15}
\loigiai{
\textbf{Đáp số:} $15$\\
$24$ phút $= 0{,}4$ giờ $\Rightarrow 6 : 0{,}4 = 15\text{ km/h}$.
}
\end{ex}

\begin{ex} %[5D12V73]
Một máy bay cất cánh lúc $8$ giờ $30$ phút và bay quãng đường $1,400\text{ km}$ với vận tốc $700\text{ km/h}$. Hỏi máy bay hạ cánh lúc mấy giờ mấy phút? Tính số phút từ $0$ giờ đến lúc hạ cánh.
\shortans{630}
\loigiai{
\textbf{Đáp số:} $630$\\
Thời gian bay: $1,400 : 700 = 2$ giờ.\\
Hạ cánh: $10$ giờ $30$ phút $= 10 \times 60 + 30 = 630$ phút.
}
\end{ex}

\begin{ex} %[5D12C73]
Một người đi xe máy từ A đến B với vận tốc $40\text{ km/h}$. Sau $45$ phút, một ô tô xuất phát từ A đuổi theo với vận tốc $55\text{ km/h}$. Hỏi sau bao nhiêu giờ ô tô đuổi kịp xe máy?
\shortans{2}
\loigiai{
\textbf{Đáp số:} $2$\\
Khoảng cách ban đầu: $40 \times 0{,}75 = 30\text{ km}$.\\
Thời gian đuổi kịp: $30 : (55 - 40) = 30 : 15 = 2\text{ giờ}$.
}
\end{ex}
